% =========================================================
% ANEXO A - MANUAL DE USUARIO (contenido)
% Este archivo se incluye desde el main, donde YA se hace:
% \appendix
% \chapter{Anexo del Manual de Usuario}
% % =========================================================
% ANEXO A - MANUAL DE USUARIO (contenido)
% Este archivo se incluye desde el main, donde YA se hace:
% \appendix
% \chapter{Anexo del Manual de Usuario}
% % =========================================================
% ANEXO A - MANUAL DE USUARIO (contenido)
% Este archivo se incluye desde el main, donde YA se hace:
% \appendix
% \chapter{Anexo del Manual de Usuario}
% % =========================================================
% ANEXO A - MANUAL DE USUARIO (contenido)
% Este archivo se incluye desde el main, donde YA se hace:
% \appendix
% \chapter{Anexo del Manual de Usuario}
% \input{.../manual_de_usuario}
% =========================================================

% =========================================================
% PARCHE ToC: evita "A.1Introducción" (sin espacio)
% Causa: numwidth del ToC demasiado pequeño con numeración tipo A.1
% Solución: aumentar el ancho reservado a la numeración en secciones
% =========================================================
\makeatletter
\renewcommand*\l@section{\@dottedtocline{1}{1.5em}{4.0em}}
\renewcommand*\l@subsection{\@dottedtocline{2}{3.8em}{4.6em}}
\renewcommand*\l@subsubsection{\@dottedtocline{3}{7.0em}{5.2em}}
\makeatother

% =========================
% INTRODUCCIÓN
% =========================

% ------------------------------------------------------------
% PRESENTACIÓN DEL MANUAL (ADMINISTRADOR Y TRABAJADOR)
% ------------------------------------------------------------

\paragraph*{}Esta aplicación ha sido desarrollada con el objetivo de digitalizar y centralizar los procesos esenciales de gestión de personal y planificación de eventos dentro de una organización. La plataforma permite registrar eventos, solicitar disponibilidad, asignar trabajadores, mantener un canal de comunicación interno y gestionar el seguimiento de pagos, ofreciendo un entorno único y accesible que reduce la dependencia de herramientas externas.


% =========================================================

% =========================================================
% PARCHE ToC: evita "A.1Introducción" (sin espacio)
% Causa: numwidth del ToC demasiado pequeño con numeración tipo A.1
% Solución: aumentar el ancho reservado a la numeración en secciones
% =========================================================
\makeatletter
\renewcommand*\l@section{\@dottedtocline{1}{1.5em}{4.0em}}
\renewcommand*\l@subsection{\@dottedtocline{2}{3.8em}{4.6em}}
\renewcommand*\l@subsubsection{\@dottedtocline{3}{7.0em}{5.2em}}
\makeatother

% =========================
% INTRODUCCIÓN
% =========================

% ------------------------------------------------------------
% PRESENTACIÓN DEL MANUAL (ADMINISTRADOR Y TRABAJADOR)
% ------------------------------------------------------------

\paragraph*{}Esta aplicación ha sido desarrollada con el objetivo de digitalizar y centralizar los procesos esenciales de gestión de personal y planificación de eventos dentro de una organización. La plataforma permite registrar eventos, solicitar disponibilidad, asignar trabajadores, mantener un canal de comunicación interno y gestionar el seguimiento de pagos, ofreciendo un entorno único y accesible que reduce la dependencia de herramientas externas.


% =========================================================

% =========================================================
% PARCHE ToC: evita "A.1Introducción" (sin espacio)
% Causa: numwidth del ToC demasiado pequeño con numeración tipo A.1
% Solución: aumentar el ancho reservado a la numeración en secciones
% =========================================================
\makeatletter
\renewcommand*\l@section{\@dottedtocline{1}{1.5em}{4.0em}}
\renewcommand*\l@subsection{\@dottedtocline{2}{3.8em}{4.6em}}
\renewcommand*\l@subsubsection{\@dottedtocline{3}{7.0em}{5.2em}}
\makeatother

% =========================
% INTRODUCCIÓN
% =========================

% ------------------------------------------------------------
% PRESENTACIÓN DEL MANUAL (ADMINISTRADOR Y TRABAJADOR)
% ------------------------------------------------------------

\paragraph*{}Esta aplicación ha sido desarrollada con el objetivo de digitalizar y centralizar los procesos esenciales de gestión de personal y planificación de eventos dentro de una organización. La plataforma permite registrar eventos, solicitar disponibilidad, asignar trabajadores, mantener un canal de comunicación interno y gestionar el seguimiento de pagos, ofreciendo un entorno único y accesible que reduce la dependencia de herramientas externas.


% =========================================================

% =========================================================
% PARCHE ToC: evita "A.1Introducción" (sin espacio)
% Causa: numwidth del ToC demasiado pequeño con numeración tipo A.1
% Solución: aumentar el ancho reservado a la numeración en secciones
% =========================================================
\makeatletter
\renewcommand*\l@section{\@dottedtocline{1}{1.5em}{4.0em}}
\renewcommand*\l@subsection{\@dottedtocline{2}{3.8em}{4.6em}}
\renewcommand*\l@subsubsection{\@dottedtocline{3}{7.0em}{5.2em}}
\makeatother

% =========================
% INTRODUCCIÓN
% =========================

% ------------------------------------------------------------
% PRESENTACIÓN DEL MANUAL (ADMINISTRADOR Y TRABAJADOR)
% ------------------------------------------------------------

\paragraph*{}Esta aplicación ha sido desarrollada con el objetivo de digitalizar y centralizar los procesos esenciales de gestión de personal y planificación de eventos dentro de una organización. La plataforma permite registrar eventos, solicitar disponibilidad, asignar trabajadores, mantener un canal de comunicación interno y gestionar el seguimiento de pagos, ofreciendo un entorno único y accesible que reduce la dependencia de herramientas externas.

